\documentclass[10pt,a4paper]{amsart}
\usepackage[margin=19mm]{geometry}
\usepackage{amsmath,amssymb,mathtools}
\usepackage[scr=rsfs,cal=euler]{mathalfa}
\usepackage{xfrac}
\usepackage[unicode,hidelinks,bookmarks=false]{hyperref}
\newcommand{\ie}{i.e.\ }
\newcommand{\cf}{cf.\ }
\newcommand{\resp}{respectively\ }
\newcommand{\Subsection}{\subsection}
\providecommand{\nobreakdash}{\nobreak\hbox{-}}
\setlength{\emergencystretch}{2em}
\multlinegap0pt
\usepackage{bm}
\usepackage{bbm}
\usepackage{dsfont}
\usepackage{xspace}
\usepackage{cancel}
\usepackage{mathtools}
\usepackage{cleveref}
\usepackage{amsthm}
\makeatletter
\@ifundefined{theorem}{\newtheorem{theorem}{Theorem}}{}
\@ifundefined{observation}{\newtheorem{observation}[theorem]{Observation}}{}
\makeatother
\title[Patterns that require distinct singular values]{Patterns that require distinct singular values}
\author{Caleb Cheung, Bryan Shader}
\date{Source: arXiv:2511.12641v1 (CC BY 4.0)}
\begin{document}
\maketitle

\noindent\emph{Source and licence.} Patterns that require distinct singular values. Version arXiv:2511.12641v1 (CC BY 4.0), distributed under the Creative Commons Attribution 4.0 International licence, as recorded in the arXiv licence metadata for this exact version and shown on the abstract page https://arxiv.org/abs/2511.12641v1. Full rights notice: licenses/arxiv-2511.12641-CC-BY-4.0.txt.\par\smallskip

\noindent\textit{Editorial scope.} Counted unit: the Theorem whose source label is \texttt{bcaterpillars} in the article (source0.tex lines 1307--1312, source label \texttt{bcaterpillars}), delivered together with the article's own complete original proof of it (source0.tex lines 1313--1420, one \texttt{proof} environment of 108 lines). No mathematics is written, repaired or re-derived: statement and proof are the article's own text line for line. Required context retained uncounted: fiedlergraphdistinct (lines 1151--1155); alternating (lines 1292--1304). Every definition, display and label of those lines is retained unchanged. Omitted from this package and not redistributed: 1--1150, 1156--1291, 1305--1306, 1421--1755. Nothing inside a retained excerpt is omitted or altered: every retained body is delivered exactly as the article's lines are written, so no omission receipt of any kind accompanies this package. Citations: the retained excerpts carry no citation command at all, so no bibliography entry of the article is transcribed and no reference is redistributed. Reference adaptation: none was needed and none was made; every reference inside the delivered unit resolves inside it, so no unresolved reference or citation is produced. Printed slips kept literal: no printed word, symbol, hypothesis or number of the counted statement or of its proof was altered. Layout and class: the article's own class is \texttt{elsarticle} with packages including graphicx, dsfont, amsmath, amssymb, amsthm, xcolor, mathabx, float, multirow, hyperref, amsbsy, algorithmic, algorithm2e, cleveref; the wrapper is the lane's pinned \texttt{amsart} editorial wrapper at 10pt on A4, which restates the theorem-like environments the retained text needs and adds the article's own macro definitions that the retained text needs (none), with extra wrapper packages (bm, bbm, dsfont, xspace, cancel, mathtools, cleveref). The delivered numbering is the unit's own. Assets: no retained span contains a figure environment, a graphics-inclusion command, a PDF-page inclusion, a table or a TikZ picture; no image file is copied into this package, the package declares no asset dependency and no image file is redistributed with it. Licence: version arXiv:2511.12641v1 of this article is distributed under the Creative Commons Attribution 4.0 International licence, as recorded in the arXiv licence metadata for this exact version and shown on the abstract page https://arxiv.org/abs/2511.12641v1. A full rights notice is delivered at licenses/arxiv-2511.12641-CC-BY-4.0.txt. Characters: every retained excerpt is pure ASCII, so the delivered engine, XeLaTeX with its default Latin Modern text font, reports no missing character.
\begin{observation}
\label{fiedlergraphdistinct}
Let $P$ by an $m\times m$ matrix of term-rank $m$ whose bigraph is connected.  Then no matrix in $\mathcal{Z}(P)$ has a multiple singular value if and only if $\mathcal{B}(P)$
is a Fiedler graph. 
\end{observation}

\begin{observation}
\label{alternating}
Let $G$ be a bipartite graph, $v$ be a vertex of $G$, and  $M$ be a matching of $G\setminus \{v\}$.
Suppose that 
\[ 
\mbox{ $i_1$---$i_2$--- $\cdots$ ---$i_{2\ell}$}
\]
is an $M$-alternating path in $G$ with $v=i_1$.
Then replacing the edges $i_{2k}$---$i_{2k+1}$ 
for $k=1,2,\ldots, \ell-1$ of $M$ by the edges 
$i_{2k-1}$---$i_{2k}$ $k=1,\ldots, \ell-1$
results in a matching of $G\setminus \{i_{2\ell}\}$ with the same number of edges as $M$. 
\end{observation}

\begin{theorem}
\label{bcaterpillars}
Let $P$ be an $m\times n$ pattern with $m\leq n$, term-rank $m$, 
and no columns of all zeros. If $P$ does not allow a matrix with a multiple singular value, then the bigraph of $P$ 
is an $m\times n$ Fiedler graph. 
\end{theorem}

\begin{proof}
   Assume no matrix in $\mathcal{Z}(P)$ 
  has a multiple singular value.  We show that $\mathcal{B}(P)$ is an $m\times n$
  Fiedler graph by establishing a sequence of structural properties of 
  $\mathcal{B}(P)$.  We indicate the end of each proof of a claim by a $\diamond$ .

\bigskip\noindent
{\bf Claim 1.} $\mathcal{B}(P)$ is connected. 
\\[12pt]
    Proof. Suppose not.  Then, after row and column permutations, 
    $P$ is  the direct sum of two nonzero patterns.  There exist matrices with these patterns with the same largest singular values. The resulting direct sum has pattern $P$ and a multiple singular value, contrary to assumption. \hfill $\diamond$ 

\bigskip\noindent
{\bf Claim 2.} Let $\alpha \subseteq \{1,\ldots, n\}$ such that 
$|\alpha|=m$ and the bipartite subgraph $B_{\alpha}$ of $P$ consisting of all row vertices 
and the column vertices in $\alpha$ has an $m$-matching.
Then each connected component of  $B_{\alpha}$
is a Fiedler graph.
\\ [12pt]
Proof. Without loss of generality, we may assume that $\alpha=\{1,\ldots, m\}$.
As $P[\{1,\ldots,m\}]$ has term-rank $m$, each of the connected components of $B_{\alpha}$ 
has a perfect matching.  If one of these components is not a Fiedler graph, then, by   \Cref{fiedlergraphdistinct} and the Direct Sum theorem, there is 
a matrix in $P[\{1,\ldots,m\}]$  with the SSVP and a multiple singular 
value. Hence by the Superpattern theorem, we are led to the contradiction that $P$ allows a matrix with a multiple singular value.
Therefore, each connected component of $B_{\alpha}$ is a Fiedler graph. \hfill $\diamond$ 

\bigskip\noindent
{\bf Claim 3}.
There exists $\alpha \subseteq \{1,\ldots, n\}$ such that 
$|\alpha|=m$ and $B_{\alpha}$ is a connected graph with a perfect matching. 
\\ [12pt]
Proof. Suppose that $B_{\alpha}$ has a perfect matching $M$ and two or more components.
We show how to find an $\alpha'$ such that $B_{\alpha}'$ has a perfect matching 
and fewer connected components. 

Without loss of generality, we may take $\alpha=\{1,\ldots, m\}$.
Claim 1 implies that there exist two connected components $B_1$ and $B_2$ of $B_{\alpha}$
and a column vertex $c$ with $c>m$ such that $c$ has a neighbor $r_1$ in $B_1$
and a neighbor $r_2$ in $B_2$.
Necessarily, there exist vertices $d_1$ and $d_2$ of $B_{\alpha}$
such that {$r_1$---$d_1$} and {$r_2$---$d_2$} are edges in $M$.




If $d_1$ is a pendant vertex of $B_{\alpha}$, then $B_{\alpha \cup \{ c\} \setminus \{ d_1\}}$ has an  $m$-matching
(namely that obtained from $M$ by replacing $r_1$---$d_1$ by $c$---$r_1$), 
and one fewer connected component.
Similarly, we are done if $d_2$ is a pendant vertex in $B_{\alpha}$. 

Now assume that neither $d_1$ nor $d_2$ is a pendant vertex in $B_{\alpha}$. If $r_1$ has degree two or more in $B_\alpha$, 
then $B_{\alpha\cup \{c\}}$ contains a subgraph that is the 
disjoint union of  a subdivided claw with $r_1$ as the center, and a matching 
of size $m-3$. The Direct Sum and Superpattern theorems now 
lead to the contradiction that $P$ allows a matrix with a multiple singular value.  Thus $r_1$, and similarly, $r_2$ are pendant vertices
in $B_{\alpha}$.  As $B_{\alpha \cup \{c\} }$ has an $m$-matching, 
there exists an $M$-alternating path starting at $c$ and containing 
$r_1$. Let $c'$ be the terminal vertex of such a 
path. By \Cref{alternating}, $B_{\{\alpha \cup \{ c\} \setminus{c'} }$
has an $m$-matching and fewer connected components than $B_{\alpha}$.
\hfill $\diamond$ 



\bigskip 
By Claims 1--3, we can assume that $B_{\{ 1,2,\ldots, m\}}$
is an $m\times m$ Fiedler graph. In particular, it is a tree
with a perfect matching $M$. Let $c>m$ be a column vertex. 

\bigskip\noindent 
{\bf Claim 4. } $c$ is a pendant vertex in $\mathcal{B}(P)$.
\\[12pt]
Proof. Suppose to the contrary that $c$ is not pendant. 
Let $r_1$ and $r_2$ be two of its neighbors.
Let $c_1$ and $c_2$ be the vertices such that $r_1$---$c_1$ and $r_2$---$r_2$' are edges of $M$. 

There exists an $M$-alternating path 
in $B_{\{1,2,\ldots, m,c\}}$ that starts at $c$, contains $r_1$ 
and ends at some pendant column vertex $c'$.  By \Cref{alternating}
$B_{\{1, \ldots, m\}} \cup \{c\} \setminus \{c'\}\}$ contains 
an $M$-matching, but also a cycle. Thus  $B_{\{1, \ldots, m\} \cup \{c\} \setminus \{c'\}\}}$ is not a Fiedler graph, contrary to Claim 1.
Therefore, $c$ is  a pendant vertex in $\mathcal{B}(P)$. \hfill $\diamond$ 


\bigskip\noindent
Let $\gamma$ be the path of $B_{\{1,\ldots, m\}}$
given in the definition of an ($m\times m$) Fiedler graph.


\bigskip\noindent
{\bf Claim 5.}  The neighbor $r$  of $c$ is not a pendant vertex
of $B_{\{1,2,\ldots, m\}}$, unless $r$ is the 
first or last vertex of the path $\gamma$. 
\\ [12pt]
Proof. 
Suppose not. Let $c'$ be the neighbor of $r$ in $B_{\{1,\ldots, n\}}$.  As $\mathcal{B}$ has a perfect matching $M$,
and $c'$ is not pendant,  $r$--$c'$ is in $M$, $c'$ has two additional neighbors $r_1$ and $r_2$ on $\gamma$  Let $r_1$---$c_1$ and $r_2$---$c_2$ be the edges of $M$ that contain $r_1$ and $r_2$.
Note that  the induced subgraph $H$ with vertices $c,r,c',r_1,r_2,c_1,c_2$
is a subdivided claw. Hence $B_{\{1,\ldots, m,c\}}$
has a subgraph that is the disjoint union of a subdivided claw 
and a $(m-3)$-matching.  As before, this contradicts that $P$ does not
allow a matrix with a multiple singular value. \hfill $\diamond$ 



\bigskip\noindent
Claims 1-5 now imply that $\mathcal{B}(P)$ is an $m\times n$ Fiedler graph.
\end{proof}

\end{document}
